%% kapa: exam=1 round=29 year=2018 period=1 form=A subject=economics subject_ko=경제학원론 seq=6 number=46
%% source: Q-Net 감정평가사 기출문제 (cst00309) artlSeq=5202808 file=extracted/1008/5202808/2018년 제29회 감정평가사 1차시험 기출문제/2018년 제29회 감정평가사 1차시험 기출문제/1교시 A형.pdf
%% conversion: kordoc 4.24.51 -> tools/gen1.py (auto); review=proofread-sonnet; flags=none
\begin{question}{46}
甲의 효용함수는 $U=\sqrt{LF}$ 이며 하루 24시간을 여가( $L$)와 노동( $24-L$)에 배분한다. 甲은 노동을 통해서만 소득을 얻으며, 소득은 모두 식품( $F$)을 구매하는 데 사용한다. 시간당 임금은 10,000원, 식품의 가격은 2,500원이다. 甲이 예산제약 하에서 효용을 극대화할 때, 여가시간과 구매하는 식품의 양은?
\begin{choices}
  \item $L=8,\ F=64$
  \item $L=10,\ F=56$
  \item $L=12,\ F=48$
  \item $L=14,\ F=40$
  \item $L=16,\ F=32$
\end{choices}
\end{question}
\answer{3}
